\subsection{Spectral Tensor Form of the Quartic Distortion Problem}

The standard quartic centrifugal distortion problem may be formulated in terms
of the compressed pair--pair tensor
\begin{equation}
\boldsymbol{\tau}=
(\tau_{aaaa},\tau_{bbbb},\tau_{cccc},\tau_{aabb},\tau_{aacc},\tau_{bbcc})^T .
\end{equation}
For the purpose of projection onto spectroscopic quartic constants, it is
convenient to introduce the Gaussian-consistent reduced matrix
\begin{equation}
\Tau'=
\begin{pmatrix}
\tau_{aaaa} & \tau_{aabb}/2 & \tau_{aacc}/2 \\
\tau_{aabb}/2 & \tau_{bbbb} & \tau_{bbcc}/2 \\
\tau_{aacc}/2 & \tau_{bbcc}/2 & \tau_{cccc}
\end{pmatrix}.
\label{eq:tauprime-matrix}
\end{equation}
Because $\Tau'$ is a real symmetric $3\times 3$ matrix, it admits the spectral
representation
\begin{equation}
\Tau' = R \,\mathrm{diag}(\lambda_1,\lambda_2,\lambda_3)\, R^T,
\label{eq:tauprime-spectral}
\end{equation}
where $R\in SO(3)$ is the eigenframe rotation matrix and
$\lambda_1,\lambda_2,\lambda_3$ are the eigenvalues of $\Tau'$.
Equation~\eqref{eq:tauprime-spectral} provides a natural tensorial
parameterization of the quartic distortion problem in terms of three spectral
invariants and three orientation parameters.

This $3+3$ description is, however, redundant with respect to the five
independent quartic Watson constants.  The reduction from six tensorial
coordinates to five spectroscopic parameters is not obtained by discarding one
angle, but rather through the existence of a one-dimensional null direction of
the linear map from $\boldsymbol{\tau}$ to Watson constants.

For a fixed reduction (\emph{e.g.}\ $A$ or $S$), the quartic projection
\begin{equation}
\boldsymbol{\tau} \longmapsto
(\Delta_J,\Delta_{JK},\Delta_K,\delta_J,\delta_K)
\end{equation}
has rank five.  Therefore the compressed quartic tensor is determined only up
to addition of a one-dimensional null vector.  Writing
\begin{equation}
\Sigma=\frac{2A-B-C}{B-C},
\end{equation}
one finds that the null direction in the compressed pair--pair basis is
\begin{equation}
\mathbf{n}=
\left(
0,\,
0,\,
0,\,
1,\,
\frac{\Sigma-1}{2},\,
-\frac{\Sigma+1}{2}
\right)^T,
\label{eq:null-vector}
\end{equation}
up to an arbitrary overall normalization.  Equivalently, in matrix form the
corresponding null tensor is
\begin{equation}
N'=
\begin{pmatrix}
0 & 1/2 & (\Sigma-1)/4 \\
1/2 & 0 & -(\Sigma+1)/4 \\
(\Sigma-1)/4 & -(\Sigma+1)/4 & 0
\end{pmatrix}.
\label{eq:null-matrix}
\end{equation}
Thus all matrices of the form
\begin{equation}
\Tau' + c\,N'
\end{equation}
lead to the same set of quartic Watson constants.

The Moore--Penrose pseudoinverse used in the tensor reconstruction algorithm
does not define the physics itself, but selects a particular representative
within this affine family.  In the compressed pair--pair basis this
pseudoinverse gauge fixing is expressed by the orthogonality condition
\begin{equation}
\tau_{aabb}
\;+\;
\frac{\Sigma-1}{2}\tau_{aacc}
\;-\;
\frac{\Sigma+1}{2}\tau_{bbcc}
=0,
\label{eq:pseudoinverse-gauge-tau}
\end{equation}
which is equivalent to
\begin{equation}
2\Tau'_{12}
\;+\;
(\Sigma-1)\Tau'_{13}
\;-\;
(\Sigma+1)\Tau'_{23}
=0.
\label{eq:pseudoinverse-gauge-tauprime}
\end{equation}
Equation~\eqref{eq:pseudoinverse-gauge-tauprime} should therefore be understood
as a gauge-fixing condition associated with the pseudoinverse reconstruction,
not as a physical admissibility condition on quartic distortion tensors.

This observation clarifies the geometric structure of the quartic problem.  The
fundamental tensorial object is the symmetric matrix $\Tau'$, whose spectral
coordinates provide a natural $3+3$ description in $R^3$.  The usual Watson
quartic constants correspond to the quotient of this tensorial description by
the one-dimensional null family generated by Eq.~\eqref{eq:null-matrix}.  In
this sense, the practical five-parameter spectroscopic description is most
naturally viewed as a reduced coordinate system on the quotient space
\begin{equation}
\mathrm{Sym}(3)\big/\mathrm{span}(N').
\end{equation}

The general tensorial character of quartic centrifugal distortion constants has
been recognized for a long time, and invariant combinations of rotational
Hamiltonian parameters have also been studied in the literature.\cite{Yamada1973,Kirschner1983}
What is important in the present context is that, for the standard
$H_{12}H_{12}$ contribution and its representation-transform problem, the
Gaussian-consistent matrix $\Tau'$ provides a particularly transparent bridge
between the full tensorial description and the practical Watson
parameterization.  To the best of our knowledge, the explicit formulation in
terms of the spectral decomposition of $\Tau'$ together with the associated
pseudoinverse gauge fixing of Eqs.~\eqref{eq:pseudoinverse-gauge-tau} and
\eqref{eq:pseudoinverse-gauge-tauprime} has not been used as the organizing
principle for quartic representation transforms in this form.
